% GENERATED FILE. Edit canonical lessons, then run npm run generate:books.
% canonical-source-hash: fnv1a64:76eb6469e460ca72
% canonical-lessons: TA-C316-pinnar, TA-C316-anraiya, TA-C316-muntaiya, TA-C316-atanpin, TA-C316-ate-camayam

\chapter{That Very Day}
\label{ch:ta-that-very-day}
\hlchaptermodality{316}

\begin{chapteropening}
\textbf{By the end of this chapter:} \emph{I can place events on a day in a story.}

Complete the last lesson of chapter 316: I can place events on a day in a story.
\end{chapteropening}

\section[\texorpdfstring{\ta{பின்னர்} (piṉṉar)}{பின்னர் (piṉṉar)}]{\ta{பின்னர்} (piṉṉar) — later on; afterwards}

\label{lesson:TA-C316-pinnar}



\begin{quote}
Before the new one: say the Tamil for at present, then the Tamil for once in a while.
\end{quote}

\subsection*{You'll want to know: \ta{பின்னர்}}
\textbf{\ta{பின்னர்}} — \emph{piṉṉar} — \textquotedblleft{}later on; afterwards\textquotedblright{}.

\begin{quote}
\ta{முதலில்} \ta{நாங்கள்} \ta{சாப்பிட்டோம்}, \ta{பின்னர்} \ta{கடைக்குப்} \ta{போனோம்}.
\end{quote}

\begin{quote}
\emph{mutalil nāṅkaḷ cāppiṭṭōm, piṉṉar kaṭaikkup pōṉōm.}
\end{quote}

\begin{quote}
First we ate, and later on we went to the shop.
\end{quote}

A more formal word for \ta{பிறகு}.

\begin{cousinweb}
\textbf{\ta{பின்னர்}} (\emph{piṉṉar}) is native Dravidian: it is built on \ta{பின்}, \textquotedblleft{}back, what comes after\textquotedblright{} (DEDR 4205).
\end{cousinweb}

\begin{grammarlens}[title={What you've built}]
One more word for telling what happened.
\end{grammarlens}

\subsection*{Your turn}
\begin{itemize}
\raggedright
  \item \emph{Say it:} \emph{piṉṉar}
  \item \emph{Say it:} \emph{piṉṉar}, once more
  \item \emph{Say it:} \emph{mutalil nāṅkaḷ cāppiṭṭōm, piṉṉar kaṭaikkup pōṉōm.} — First we ate, and later on we went to the shop.
  \item \emph{Recall:} say the Tamil for at present, then the Tamil for once in a while, then say \emph{piṉṉar} again
\end{itemize}

\begin{tcolorbox}[breakable,colback=teal!4,colframe=teal!35!black,title={Before you move on}]
What does \ta{பின்னர்} mean? (\textquotedblleft{}Later on; afterwards\textquotedblright{}.) Say it once more.
\end{tcolorbox}
\section[\texorpdfstring{\ta{அன்றைய} (aṉṟaiya)}{அன்றைய (aṉṟaiya)}]{\ta{அன்றைய} (aṉṟaiya) — of that day; that day's}

\label{lesson:TA-C316-anraiya}



\begin{quote}
Before the new one: say the Tamil for once in a while, then the Tamil for later on.
\end{quote}

\subsection*{You'll want to know: \ta{அன்றைய}}
\textbf{\ta{அன்றைய}} — \emph{aṉṟaiya} — \textquotedblleft{}of that day; that day's\textquotedblright{}.

\begin{quote}
\ta{அன்றைய} \ta{நாளிதழில்} \ta{அந்தச்} \ta{செய்தி} \ta{வந்தது}.
\end{quote}

\begin{quote}
\emph{aṉṟaiya nāḷitaḻil antac ceyti vantatu.}
\end{quote}

\begin{quote}
That news appeared in that day's newspaper.
\end{quote}

Like \ta{இன்றைய}, today's, and \ta{நேற்றைய}, yesterday's; only before a noun.

\begin{cousinweb}
\textbf{\ta{அன்றைய}} (\emph{aṉṟaiya}) is native Dravidian: DEDR 1 lists \ta{அன்று} and \ta{அன்றை}, \textquotedblleft{}that day\textquotedblright{}, built on the far-pointing \ta{அ}-.
\end{cousinweb}

\begin{grammarlens}[title={What you've built}]
One more word for telling what happened.
\end{grammarlens}

\subsection*{Your turn}
\begin{itemize}
\raggedright
  \item \emph{Say it:} \emph{aṉṟaiya}
  \item \emph{Say it:} \emph{aṉṟaiya}, once more
  \item \emph{Say it:} \emph{aṉṟaiya nāḷitaḻil antac ceyti vantatu.} — That news appeared in that day's newspaper.
  \item \emph{Recall:} say the Tamil for once in a while, then the Tamil for later on, then say \emph{aṉṟaiya} again
\end{itemize}

\begin{tcolorbox}[breakable,colback=teal!4,colframe=teal!35!black,title={Before you move on}]
What does \ta{அன்றைய} mean? (\textquotedblleft{}Of that day; that day's\textquotedblright{}.) Say it once more.
\end{tcolorbox}
\section[\texorpdfstring{\ta{முந்தைய} (muntaiya)}{முந்தைய (muntaiya)}]{\ta{முந்தைய} (muntaiya) — previous; earlier}

\label{lesson:TA-C316-muntaiya}



\begin{quote}
Before the new one: say the Tamil for later on, then the Tamil for of that day.
\end{quote}

\subsection*{You'll want to know: \ta{முந்தைய}}
\textbf{\ta{முந்தைய}} — \emph{muntaiya} — \textquotedblleft{}previous; earlier\textquotedblright{}.

\begin{quote}
\ta{முந்தைய} \ta{நாள்} \ta{இரவு} \ta{நல்ல} \ta{மழை} \ta{பெய்தது}.
\end{quote}

\begin{quote}
\emph{muntaiya nāḷ iravu nalla maḻai peytatu.}
\end{quote}

\begin{quote}
The night before, it rained heavily.
\end{quote}

Only before a noun: \ta{முந்தைய} \ta{வாரம்}, the previous week.

\begin{cousinweb}
\textbf{\ta{முந்தைய}} (\emph{muntaiya}) is native Dravidian, from \ta{முந்து}, \textquotedblleft{}to come in front, be prior in time\textquotedblright{} (DEDR 5020).
\end{cousinweb}

\begin{grammarlens}[title={What you've built}]
One more word for telling what happened.
\end{grammarlens}

\subsection*{Your turn}
\begin{itemize}
\raggedright
  \item \emph{Say it:} \emph{muntaiya}
  \item \emph{Say it:} \emph{muntaiya}, once more
  \item \emph{Say it:} \emph{muntaiya nāḷ iravu nalla maḻai peytatu.} — The night before, it rained heavily.
  \item \emph{Recall:} say the Tamil for later on, then the Tamil for of that day, then say \emph{muntaiya} again
\end{itemize}

\begin{tcolorbox}[breakable,colback=teal!4,colframe=teal!35!black,title={Before you move on}]
What does \ta{முந்தைய} mean? (\textquotedblleft{}Previous; earlier\textquotedblright{}.) Say it once more.
\end{tcolorbox}
\section[\texorpdfstring{\ta{அதன்பின்} (ataṉpiṉ)}{அதன்பின் (ataṉpiṉ)}]{\ta{அதன்பின்} (ataṉpiṉ) — after that}

\label{lesson:TA-C316-atanpin}



\begin{quote}
Before the new one: say the Tamil for of that day, then the Tamil for previous.
\end{quote}

\subsection*{You'll want to know: \ta{அதன்பின்}}
\textbf{\ta{அதன்பின்}} — \emph{ataṉpiṉ} — \textquotedblleft{}after that\textquotedblright{}.

\begin{quote}
\ta{அதன்பின்} \ta{அவனை} \ta{யாரும்} \ta{பார்க்கவில்லை}.
\end{quote}

\begin{quote}
\emph{ataṉpiṉ avaṉai yārum pārkkavillai.}
\end{quote}

\begin{quote}
After that, nobody saw him.
\end{quote}

Like \ta{அதன்} \ta{பிறகு}; written as one word or two.

\begin{cousinweb}
\textbf{\ta{அதன்பின்}} (\emph{ataṉpiṉ}) is built from \ta{அதன்}, \textquotedblleft{}of that\textquotedblright{}, and native \ta{பின்}, \textquotedblleft{}back, what comes after\textquotedblright{} (DEDR 4205).
\end{cousinweb}

\begin{grammarlens}[title={What you've built}]
One more word for telling what happened.
\end{grammarlens}

\subsection*{Your turn}
\begin{itemize}
\raggedright
  \item \emph{Say it:} \emph{ataṉpiṉ}
  \item \emph{Say it:} \emph{ataṉpiṉ}, once more
  \item \emph{Say it:} \emph{ataṉpiṉ avaṉai yārum pārkkavillai.} — After that, nobody saw him.
  \item \emph{Recall:} say the Tamil for of that day, then the Tamil for previous, then say \emph{ataṉpiṉ} again
\end{itemize}

\begin{tcolorbox}[breakable,colback=teal!4,colframe=teal!35!black,title={Before you move on}]
What does \ta{அதன்பின்} mean? (\textquotedblleft{}After that\textquotedblright{}.) Say it once more.
\end{tcolorbox}
\section[\texorpdfstring{\ta{அதே} \ta{சமயம்} (atē camayam)}{அதே சமயம் (atē camayam)}]{\ta{அதே} \ta{சமயம்} (atē camayam) — at the same time}

\label{lesson:TA-C316-ate-camayam}



\begin{quote}
Before the new one: say the Tamil for previous, then the Tamil for after that.
\end{quote}

\subsection*{You'll want to know: \ta{அதே} \ta{சமயம்}}
\textbf{\ta{அதே} \ta{சமயம்}} — \emph{atē camayam} — \textquotedblleft{}at the same time\textquotedblright{}.

\begin{quote}
\ta{அதே} \ta{சமயம்}, \ta{வெளியே} \ta{ஒரு} \ta{சத்தம்} \ta{கேட்டது}.
\end{quote}

\begin{quote}
\emph{atē camayam, veḷiyē oru cattam kēṭṭatu.}
\end{quote}

\begin{quote}
At the same time, a noise was heard outside.
\end{quote}

\ta{அதே}, that same + \ta{சமயம்}, time.

\begin{cousinweb}
\textbf{\ta{அதே} \ta{சமயம்}} (\emph{atē camayam}) joins \ta{அதே}, \textquotedblleft{}that same\textquotedblright{}, with \ta{சமயம்}, \textquotedblleft{}time, occasion\textquotedblright{}, borrowed from Sanskrit \emph{samaya}, \textquotedblleft{}meeting, time\textquotedblright{} (CDIAL 13185).
\end{cousinweb}

\begin{grammarlens}[title={What you've built}]
One more word for telling what happened.
\end{grammarlens}

\subsection*{Your turn}
\begin{itemize}
\raggedright
  \item \emph{Say it:} \emph{atē camayam}
  \item \emph{Say it:} \emph{atē camayam}, once more
  \item \emph{Say it:} \emph{atē camayam, veḷiyē oru cattam kēṭṭatu.} — At the same time, a noise was heard outside.
  \item \emph{Recall:} say the Tamil for previous, then the Tamil for after that, then say \emph{atē camayam} again
\end{itemize}

\begin{tcolorbox}[breakable,colback=teal!4,colframe=teal!35!black,title={Before you move on}]
What does \ta{அதே} \ta{சமயம்} mean? (\textquotedblleft{}At the same time\textquotedblright{}.) Say it once more.
\end{tcolorbox}
